\documentclass[11pt]{article}
\usepackage[margin=1in]{geometry}
\usepackage{hyperref}

\begin{document}

\title{Summary: ResQ-AI Framework}
\author{ResQ-AI Team}
\date{}
\maketitle

\section*{Overview}
The ResQ-AI project proposes an uncertainty-aware multimodal framework aimed at addressing the critical challenges of urban flood prediction and emergency resource allocation. 

\section*{Key Components}
\begin{itemize}
    \item \textbf{Multimodal Prediction:} Integrates SAR, Optical, DEM, and rainfall data using a novel architecture with FiLM and cross-attention.
    \item \textbf{Uncertainty Quantification (UQ):} Benchmarks Deep Ensembles, MC Dropout, EDL, and Conformal Prediction to yield calibrated uncertainties.
    \item \textbf{Impact Assessment:} Utilizes Monte Carlo sampling to propagate pixel-level uncertainties to population and infrastructure exposure metrics.
    \item \textbf{Decision-Making:} Solves a chance-constrained stochastic optimization problem for resource allocation, using the generated impact distributions.
\end{itemize}

\section*{Significance}
By closing the loop between Earth Observation deep learning and Operations Research, ResQ-AI demonstrates that calibrated predictive uncertainty can significantly reduce unmet demand during disaster response, saving lives and resources compared to traditional deterministic pipelines.

\end{document}
